\begin{table}[!htbp]
\centering\small
\caption{Risk Indifference and the Duration--Effort Comparison}
\label{tab:r6_frontier}
\begin{tabular}{lrrrr}
\toprule
$k$ / restriction & Opposite-sign bracket & $A_+-A_-$ & $C_+-C_-$ & Ratio \\
\midrule
0.5 & [0.34687805, 0.34689331] & 0.00230304 & 4.6285e-05 & 0.0200973 \\
2 & [0.37007141, 0.37008667] & 0.00230304 & 2.12312e-05 & 0.00921875 \\
8 & [0.48182678, 0.48184204] & 0.00230304 & 7.90303e-05 & 0.0343156 \\
No adjustment & [0.55891418, 0.55892944] & 0.00230304 & 0 & 0 \\
\bottomrule
\end{tabular}
\par\smallskip
\begin{minipage}{0.98\linewidth}\footnotesize
Notes: Each bracket has opposite risk-advantage signs at its endpoints and width at most $2\times10^{-5}$. It proves existence of a crossing, not global uniqueness. Features and their ratio are evaluated at the midpoint. The ratio predicts a regular local branch slope only under the stated nonzero-duration and differentiability conditions. A zero-adjustment policy has zero adjustment effort. Bracket endpoints and their signed values are stored, rather than inferred from rounded table entries.
\end{minipage}
\end{table}
